\paragraph*{04.03.02.05 Das Projekt mit den Entscheidungs- und Berichterstattungsstrukturen sowie den Qualitätsanforderungen der Organisation in Einklang bringen}
\kcilabel{para:04.03.02.05_EntscheidungsUndBerichtsstrukturenUndQualitätsanforderungen}{04.03.02.05}
\label{para:04.03.02.05_EntscheidungsUndBerichtsstrukturenUndQualitätsanforderungen}

%Task
\par Aufgrund seiner Natur mit basaler Relevanz für die \gls{gob}, der Beteiligung von sieben externen Partnern und aller Konzernunternehmen stellte das Programm \gls{r2h} hohe Anforderungen an Entscheidungs- und Berichtsstrukturen samt Qualitätsanforderungen. Als \gls{wsl} stellte ich nahtlose Datenbelieferung im \gls{s4} sicher und etablierte adäquate Berichtsstrukturen und Qualitätsanforderungen.

%Action
\par Da mein Counterpart des \gls{gu} im \gls{ws} \gls{sst} die definierten Anforderungen an wöchentliche Statusberichte nicht erfüllte, definierte und implementierte ich Berichtsstrukturen und Qualitätsanforderungen eigenständig (vgl. \autoref{fig:R2HStatusSSTJuni2025}). Gestaltete sie \gls{gob}-konform und entscheidungsrelevant. Implementierte Qualitätssicherungsprozess für gelieferte \gls{ricefw}-Objekte zur Zuverlässigkeit der Berichte (vgl.\autoref{fig:TestpläneSST}).

%Result
\par Statusberichte entsprechen \gls{pl}-Anforderungen und liefern relevante Informationen in benötigter Tiefe, Qualität und adressatengerecht. Qualitätssicherung verbesserte Qualität und Aktualität, erleichterte Entscheidungsfindung auf Managementebene. Nachweis: Fähigkeit, Berichtsstrukturen und Qualitätsanforderungen eigenständig zu definieren und zu implementieren.

%Reflect
\par Rückblickend hätte ich Anforderungen klarer definiert und kommuniziert, um Missverständnisse zu vermeiden. Regelmäßige Feedbackschleifen mit Stakeholdern wären hilfreich gewesen. Zukünftig kläre ich Standards frühzeitig und arbeite eng mit Stakeholdern zusammen, um optimale Entscheidungsunterstützung zu gewährleisten.
